\section{\label{sec:intro}Introduction}

What governs equilibrium? The canonical answer is the thermal Gibbs state, where the environment is characterised by a single parameter: the temperature. Yet, this canonical answer is puzzling. It implies that neither the specific character of the environment nor the form of its coupling to the system alters the system's equilibrium. This of course is a consequence of the assumptions used to derive the answer - the weak-coupling or Markovian regimes open-systems theory has historically relied upon \cite{breuerTheoryOpenQuantum2002}. At strong coupling however, the reduced state inherits explicit temperature dependence and interaction-induced operator content that definitionally cannot be captured by these assumptions \cite{hanggiFiniteQuantumDissipation2008,ingoldSpecificHeatAnomalies2009}. These issues have motivated a broad investigation into strong-coupling thermodynamics, including exact fluctuation relations~\cite{campisiFluctuationTheoremArbitrary2009,jarzynskiNonequilibriumWorkTheorem2004,gooldNonequilibriumThermodynamics2019}, consistent heat and work definitions~\cite{espositoNatureHeatStrongly2015,rivasStrongCouplingThermodynamics2020,lacerdaQuantumThermodynamicsFast2023}, and operational measurability~\cite{strasbergMeasurabilityNonequilibriumThermodynamics2020,millerEntropyProduction2017,binderOperationalThermodynamics2019}.

While formally one can write the reduced state in canonical exponential form, the effective Hamiltonian governing this state is generally not the bare Hamiltonian of the system. Instead, it must account for the non-trivial influence of the environment in a manner that goes beyond the specification of temperature. If an effective operator capturing this influence can be constructed, the standard description of equilibrium holds generically across all coupling strengths. \emph{The Hamiltonian of mean force} (HMF) is the operator that reproduces this reduced equilibrium object as a Gibbs-like state, and provides the standard starting point for strong-coupling thermodynamics and mean-force Gibbs-state constructions
\cite{campisiFluctuationTheoremArbitrary2009,talknerColloquiumStatisticalMechanics2020,trushechkinOpenQuantumSystem2022,seifertFirstSecondLaw2016,millerHamiltonianMeanForce2018}.

The ultimate target of thermodynamic descriptions is therefore unambiguous, and in fact a formal definition of HMF is equally unambiguous - it is simply the logarithm of the total state with the environment traced out. It is the \emph{representation} of this operator that remains a challenge. Is the HMF a local operator? Does it admit a compact description in terms of few-body interactions? To address this, a veritable zoo of approximations - polaron transforms~\cite{jangNonequilibriumTheoryElectronic2008}, weak-coupling expansions~\cite{cresserWeakUltrastrongCoupling2021a}, reaction-coordinate mappings~\cite{duGeneralizedHamiltonianMeanforce2025a,nazirReactionCoordinateMapping2018} - have proliferated strong-coupling thermodynamics. These approximations often disagree in intermediate regimes however, and while powerful numerical techniques such as hierarchical equations of motion (HEOM) or path-integral Monte Carlo exist~\cite{moixEquilibriumreducedDensityMatrix2012,chenRigorousStochasticMatrix2014,makriExploitingClassicalDecoherence2014,tanimuraReducedHierarchicalEquations2014,songCalculationCorrelatedInitial2015}, they typically do not yield compact operator forms for the HMF.

Ultimately, the source of these difficulties lies in the trace operation, which mixes the bath and system degrees of freedom in a way that is difficult to invert. For this reason, a general framework for representing the HMF has proved elusive. The purpose of this paper is to lay the groundwork for just such a framework. To do so, we construct a representation of the equilibrium density matrix that disentangles the statistical properties of the environment from the algebraic response of the system. We term this object the \emph{quenched density}. By employing an imaginary-time influence functional~\cite{feynmanTheoryGeneralQuantum1963a,caldeiraQuantumTunnellingDissipative1983a,grabertQuantumBrownianMotion1988} and its Hubbard-Stratonovich decoupling~\cite{hubbardCalculationPartitionFunctions1959a,stratonovich1957QDistro,stockburgerExactNumberRepresentation2002}, we are able to rewrite the reduced equilibrium density not as a partial trace, but as a classical average over an ensemble of ``quenched'' imaginary-time propagators.
\begin{equation}
    \bar{\rho}_S(\beta) = \mathbb{E}_{\xi}\left[ \hat{\mathcal{T}} e^{-\int_0^\beta d\tau H_{\mathrm{eff}}[\xi(\tau)]} \right].
\end{equation}
In this picture, the bath is replaced by a stochastic field $\xi(\tau)$~\cite{stockburgerSimulatingSpinbosonDynamics2004,mccaulPartitionfreeApproachOpen2017c,wiedmannNonMarkovianDynamicsQuantum2020}, and the system evolves under a time-dependent effective Hamiltonian for each realisation of the noise.

This formulation effectively sets the system in a frame where it can be expressed in terms of classical trajectories for the bath degrees of freedom. The Hubbard-Stratonovich transformation then allows us to exploit this to represent the bath information locally at the operator level. The framework thus provides a rigorous bridge between the nonlocal path-integral description (where the bath is exact but the system operator structure is opaque) and the local operator description required for thermodynamics.

We apply this framework to establishing a rigorous benchmark: the linearly coupled Gaussian bath in the commuting (pure-dephasing) sector~\cite{clerkIntroductionQuantumNoise2010,breuerTheoryOpenQuantum2002}. In this limit, relevant for quantum nondemolition measurements and qubit dephasing~\cite{makhlinQuantumStateEngineering2001,palmaQuantumComputersDissipation1996}, the technical obstruction of imaginary time ordering vanishes, and the quenched density's average can be performed analytically. We use this recover the known result that the bath induces a static potential renormalisation $-\lambda f^2$ (the reorganisation energy)~\cite{campisiTalknerHanggi2009Solvable}. While this specific result is standard, its derivation within the quenched density framework serves two critical purposes. First, it demonstrates how the framework naturally recovers thermodynamic identities from path-integral primitives. Second, and most importantly, it provides a foundation for the generalisation of the framework to both non-commuting couplings and anharmonic environments.

The paper is structured as follows. In Sec.~\ref{sec:model} we define the total Hamiltonian and reduced equilibrium objects. Sec.~\ref{sec:quenched} develops the core of the framework: the Quenched Density representation of the imaginary-time propagator. Sec.~\ref{sec:influence_to_hmf} connects this to the influence functional and derives the exact commuting Gaussian benchmark. Sec.~\ref{sec:numerical} presents the numerical validation, confirming that the Quenched Density sampling, analytic prediction, and full finite-bath trace-out agree to machine precision. Sec.~\ref{sec:discussion} discusses how this Gaussian framework serves as the foundation for the future extension of the theory to non-Gaussian and non-commuting environments.
